\chapter{Conclusões e Trabalhos Futuros} \label{chap:conclusions}

	\section{Síntese dos resultados}

		\par Esta tese investigou a autenticação biométrica combinando fala fonada
		e fala imaginada, comparando extração manual de características
		(\ac{DTWPT} com agrupamento energético) e extração automatizada
		(\textit{autoencoders}), com seleção conduzida pela engenharia
		paraconsistente de características e classificação por uma rede neural de
		pulso residual profunda (\ac{DSNN}). O experimento foi organizado em duas
		fases: 208 execuções de \textit{ranking} na Fase~00 e 32 execuções de
		autenticação na Fase~01, cada uma repetida três vezes.

		\par \textbf{A extração manual venceu a seleção paraconsistente para os
		dois sinais}, com a mesma configuração --- \textit{wavelet} Haar, escala
		linear (\ac{LFCC}), Categoria~1 --- em voz e EEG. Nenhuma das 24 variantes
		de \textit{autoencoder} se aproximou do melhor extrator manual
		(\autoref{sec:respostasQuestoes}). A convergência dos dois sinais sobre a
		mesma configuração é, em si, um resultado: a \textit{wavelet} mais simples
		do conjunto avaliado, com o particionamento espectral mais simples,
		superou 22 famílias \textit{Daubechies} e dois esquemas cocleares.

		\par \textbf{O desempenho de autenticação, porém, permaneceu próximo do
		acaso.} A melhor das 32 configurações da Fase~01 --- fusão precoce,
		dependente de texto, $k$-fold plano, características brutas --- atingiu
		$\text{EER} = 0{,}4534$ e $\text{AUC} = 0{,}5452$, contra $0{,}5$ e $0{,}5$
		de um classificador aleatório. Nenhuma configuração ficou
		substancialmente abaixo desse patamar. Trata-se, portanto, de um
		\textbf{resultado negativo}: nas condições avaliadas, o sistema proposto
		não sustenta autenticação prática, e as conclusões desta tese devem ser
		lidas sob essa limitação central.

		\par Registre-se, ainda assim, o que o experimento indicou dentro desse
		regime: o EEG e as fusões superaram a voz isolada entre as configurações
		mais bem colocadas, o que é consistente com a motivação original de
		complementar a fala degradada com sinal cerebral; e a padronização das
		características piorou sistematicamente o \ac{EER} enquanto elevava o
		\ac{AUC}, divergência que sugere um problema de calibração de limiar e
		não de poder discriminativo.

	\section{Contribuições}

		\par Além dos resultados experimentais, três contribuições de natureza
		metodológica emergiram do trabalho, todas motivadas por defeitos
		detectados durante a própria execução:

		\begin{enumerate}
			\item \textbf{A métrica $D_{\text{penalizado}}$}
			(\autoref{sec:dPenalizado}). O grau de verdade $D_{1,0}$, usado
			isoladamente, premia extratores degenerados: uma rede que emite o mesmo
			vetor para toda entrada atinge $\alpha=\beta=1$ e pontua melhor que
			extratores genuinamente separáveis. A correção proposta soma um termo
			$\lambda|G_2|$ proporcional ao grau de contradição, com
			$\lambda = 2-\sqrt{2}$ escolhido de modo que os três vértices
			degenerados sejam penalizados igualmente. A métrica é função apenas de
			$\alpha$ e $\beta$, e pôde ser aplicada retroativamente sem reexecutar
			experimento algum.

			\item \textbf{O diagnóstico do piso de ruído das codificações
			temporais} (\autoref{sec:protocoloDuasFases}). A ordem observada entre
			as codificações --- \texttt{direct} $\succ$ \texttt{latency} $\succ$
			\texttt{poisson} --- coincide exatamente com a ordem prevista pelo
			ruído que cada codificação injeta \emph{antes} de qualquer aprendizado,
			calculável em fechado ($\sigma_{\text{poisson}} = 0{,}125$ para
			$T=16$). Como $\alpha$ é definido a partir da amplitude intraclasse, ele
			é estruturalmente incapaz de separar ``o codificador não aprendeu'' de
			``esta codificação tem piso estatístico irredutível''. A conclusão
			metodológica é que a comparação entre codificações temporais sob este
			critério exige controlar $T$, e não apenas medir.

			\item \textbf{A aplicação da engenharia paraconsistente à seleção de
			extratores}, com um protocolo em duas fases que separa a escolha da
			representação da medição de desempenho, evitando que a seleção do
			extrator seja contaminada pela capacidade do classificador.
		\end{enumerate}

	\section{Limitações}

		\par A limitação dominante é o \textbf{tamanho da base}: 15 locutores e
		1974 amostras. Sob protocolo de verificação, com partições disjuntas por
		locutor, cada dobra treina com um número de identidades muito aquém do
		usual em biometria de voz, e o risco de o classificador não dispor de sinal
		suficiente para generalizar é alto --- hipótese compatível com o
		desempenho próximo do acaso observado.

		\par Três limitações de desenho devem acompanhar qualquer leitura dos
		resultados:

		\begin{itemize}
			\item as três categorias de representação \textbf{não} foram
			comparadas através de um classificador comum, apenas pelo critério
			paraconsistente, de modo que a Questão~3 não foi respondida na forma em
			que fora originalmente formulada (\autoref{sec:respostasQuestoes});
			\item os \textit{autoencoders} apresentaram $\bar\beta$ acima de
			$0{,}97$ em todas as configurações, indicando latentes próximos do
			colapso; não se pode distinguir, com os dados disponíveis, ``pior
			extrator'' de ``treino insuficiente'';
			\item a Questão~4 (dependência de texto) foi avaliada sobre um sistema
			que opera no nível do acaso, o que torna a ausência de diferença entre
			os modos pouco informativa.
		\end{itemize}

	\section{Trabalhos futuros}

		\par O encaminhamento mais direto decorre das limitações acima:

		\begin{enumerate}
			\item \textbf{Ampliar a base}, em número de locutores antes que em
			amostras por locutor: sob protocolo de verificação, é a contagem de
			identidades que limita o que o classificador pode aprender a
			discriminar.

			\item \textbf{Testar o piso de ruído em $T=64$.} Como o piso decresce
			com $1/T$, quadruplicar o número de passos reduz
			$\sigma_{\text{poisson}}$ a $0{,}0625$. Se a desvantagem da codificação
			\texttt{poisson} encolher proporcionalmente, ela era artefato de
			medição; se persistir, é propriedade da codificação. O teste decide uma
			questão que a Fase~00 deixou explicitamente em aberto.

			\item \textbf{Comparar as três categorias através do classificador},
			estendendo a Fase~01 para além da combinação vencedora, de modo a
			responder à Questão~3 na forma originalmente pretendida.

			\item \textbf{Investigar a divergência EER--AUC sob padronização},
			verificando a hipótese de calibração de limiar --- se confirmada, a
			padronização passaria de prejudicial a benéfica mediante ajuste do
			limiar de decisão.

			\item \textbf{Variar orçamento de treino e capacidade dos
			\textit{autoencoders}}, para separar ``pior extrator'' de ``treino
			insuficiente'' e permitir uma comparação justa com a extração manual.
		\end{enumerate}
